\authorchapter{Projektbeschreibung und Anforderungen}{Dražen Plavšić, Ereny Micail}
\label{ch:projektbeschreibung}

\authorsection{Was ist EffSync?}{Dražen Plavšić}

\textbf{EffSync} (von \textit{Efficient Synchronization}) ist ein webbasiertes
Task-Management-System, das Teams und Einzelpersonen dabei unterstützt, Projekte
effizienter zu planen, zu organisieren und zu verfolgen. Die Anwendung kombiniert
klassische Projektmanagement-Funktionen mit modernen Kommunikationsfeatures und
bietet dabei eine übersichtliche, benutzerfreundliche Oberfläche.

Der Name \textit{EffSync} steht für die zwei Kernprinzipien der Anwendung:

\begin{itemize}
    \item \textbf{Effizienz} -- Aufgaben schnell erfassen, priorisieren
          und abschließen
    \item \textbf{Synchronisation} -- Teams in Echtzeit auf dem gleichen
          Stand halten
\end{itemize}

Das Projekt wurde im Rahmen der Diplomarbeit an der \textbf{BULME Graz}
von \textbf{Drazen Plavsic} und \textbf{Ereny Micail} der Klasse 3/4AKIFT
unter der Betreuung von \textbf{Simon Gunacker} entwickelt. Als
Versionsverwaltung wurde GitHub eingesetzt, wobei eine strukturierte
Git-Flow-Strategie mit Feature-Branches, einem Development-Branch und
einem stabilen Main-Branch verfolgt wurde.

\authorsection{Ausgangslage}{Dražen Plavšić, Ereny Micail}

\subsection{Ausgangslage -- Drazen Plavsic}

Die Ausgangslage basiert auf verfügbaren Schulressourcen wie Räumlichkeiten
und Laptops sowie den im Unterricht erworbenen Kenntnissen in
Webentwicklung, Datenbanken und Softwarearchitektur. Ziel war es, diese
Kenntnisse in einem praxisnahen Projekt anzuwenden und zu vertiefen.
Der Fokus lag dabei auf der Backend-Entwicklung mit NestJS, dem Design
der Datenbankstruktur mit PostgreSQL und Prisma ORM sowie der Umsetzung
eines umfassenden Sicherheitskonzepts.

\subsection{Ausgangslage -- Ereny Micail}

Die Ausgangslage umfasst ebenfalls die verfügbaren Schulressourcen sowie
die im Unterricht erworbenen Kenntnisse in Frontend-Entwicklung und
UI/UX-Design. Der Schwerpunkt lag auf der Entwicklung einer intuitiven
Benutzeroberfläche mit Vue.js, der Implementierung des Dashboard-Metriken-Systems
sowie der serverseitigen Rollen- und Rechteprüfung inklusive Audit-Log.

\authorsection{Zielgruppen}{Dražen Plavšić, Ereny Micail}

EffSync wurde für zwei klar definierte Zielgruppen konzipiert:

\subsection{Zielgruppe 1: Schulen (Schüler und Lehrer)}

Die erste Zielgruppe umfasst Schülerinnen, Schüler und Lehrpersonen an
Höheren Technischen Lehranstalten und anderen Bildungseinrichtungen.
Für diese Gruppe steht eine \textbf{vereinfachte Version} zur Verfügung,
die folgende Eigenschaften aufweist:

\begin{itemize}
    \item Verwaltung von Klassen, Fächern, Stundenplänen und Schulräumen
    \item Aufgabenverwaltung für Schulprojekte und Hausübungen
    \item Rollenkonzept mit Superuser, Admin, Teacher und Student
    \item Einfache Installation und Inbetriebnahme
    \item Live-Chat für Klassen- und Projektgruppen
\end{itemize}

\subsection{Zielgruppe 2: Unternehmen}

Die zweite Zielgruppe sind kleine bis mittelgroße Unternehmen, die eine
professionelle Projektmanagement-Lösung benötigen. Für diese Gruppe
steht eine \textbf{erweiterte Version} zur Verfügung:

\begin{itemize}
    \item PostgreSQL-Datenbank für hohe Performance und Skalierbarkeit
    \item Abbildung der Unternehmensstruktur (Abteilungen, Teams, Rollen)
    \item Erweitertes Sicherheitskonzept inkl.\ \gls{2FA} und Audit-Log
    \item Persistente Speicherung von Chatverläufen und Dateien
    \item Detaillierte Projektmetriken, Trends und Reporting
    \item \gls{DSGVO}-konforme Datenhaltung und Löschkonzepte
\end{itemize}

\authorsection{Funktionsübersicht}{Dražen Plavšić, Ereny Micail}

Tabelle~\ref{tab:features} gibt einen Überblick über die implementierten
Features von EffSync:

\begin{table}[H]
    \centering
    \caption{Funktionsübersicht EffSync}
    \begin{tabularx}{\textwidth}{lX}
        \toprule
        \textbf{Feature} & \textbf{Beschreibung} \\
        \midrule
        Aufgabenverwaltung   & Tasks \& Subtasks erstellen, bearbeiten, löschen \\
        Kalenderansicht      & Aufgaben nach Fälligkeitsdatum visualisieren \\
        Dashboard            & Metriken, Statusverteilung, Trends (wöchentlich/monatlich) \\
        Live-Chat            & Echtzeit-Kommunikation via WebSockets (Socket.io) \\
        Globale Suche        & Suche über Tasks, User, Projekte \\
        Benachrichtigungen   & In-App-Benachrichtigungen \\
        Profilverwaltung     & Benutzerprofile anpassen \\
        Google OAuth         & Anmeldung via Google-Konto \\
        RBAC                 & Rollenbasierte Zugriffskontrolle \\
        Audit-Log            & Nachverfolgung kritischer Operationen \\
        Mitgliederverwaltung & Explizite Task-Member-Verwaltung \\
        Klassenstruktur      & Klassen, Fächer, Stundenpläne \\
        Firmenstruktur       & Abteilungen und Teams abbilden \\
        Chatverlauf          & Persistente Speicherung (DSGVO) \\
        2FA                  & Zwei-Faktor-Authentifizierung \\
        Backup \& Restore    & Datensicherungsstrategie \\
        \bottomrule
    \end{tabularx}
    \label{tab:features}
\end{table}

\authorsection{Arbeitspaketplan (Work Breakdown Structure)}{Dražen Plavšić, Ereny Micail}
\label{sec:wbs}

Der Arbeitspaketplan gliedert das Projekt in sieben Phasen, die
chronologisch von Oktober 2025 bis April 2026 durchgeführt wurden.

\subsection{Phase 1: Kickoff \& Analyse (Oktober 2025)}

\begin{table}[H]
    \centering
    \caption{Arbeitspakete Phase 1}
    \begin{tabularx}{\textwidth}{llX}
        \toprule
        \textbf{AP-ID} & \textbf{Verantwortlich} & \textbf{Arbeitspaket} \\
        \midrule
        AP 1.1 & Gemeinsam & Projektidee definieren \& Scope festlegen \\
        AP 1.2 & Gemeinsam & Zielgruppenanalyse (Schule vs.\ Unternehmen) \\
        AP 1.3 & Gemeinsam & Technologie-Stack evaluieren \& entscheiden \\
        AP 1.4 & Gemeinsam & Git-Repository aufsetzen \& Git-Flow definieren \\
        AP 1.5 & Gemeinsam & Meilensteinplanung \& Zeitplan erstellen \\
        \bottomrule
    \end{tabularx}
    \label{tab:ap1}
\end{table}

\subsection{Phase 2: Grundarchitektur (November 2025)}

\begin{table}[H]
    \centering
    \caption{Arbeitspakete Phase 2}
    \begin{tabularx}{\textwidth}{llX}
        \toprule
        \textbf{AP-ID} & \textbf{Verantwortlich} & \textbf{Arbeitspaket} \\
        \midrule
        AP 2.1 & Drazen & Backend-Grundstruktur (NestJS) aufsetzen \\
        AP 2.2 & Drazen & Datenbankschema entwerfen (\gls{ERD}, Prisma Schema) \\
        AP 2.3 & Ereny  & Frontend-Grundstruktur (Vue.js, Vite) aufsetzen \\
        AP 2.4 & Ereny  & UI-Komponentenbibliothek \& Design-System \\
        AP 2.5 & Gemeinsam & API-Schnittstellen definieren \\
        \bottomrule
    \end{tabularx}
    \label{tab:ap2}
\end{table}

\subsection{Phase 3: Authentifizierung \& Benutzerverwaltung
  (Dezember 2025)}

\begin{table}[H]
    \centering
    \caption{Arbeitspakete Phase 3}
    \begin{tabularx}{\textwidth}{llX}
        \toprule
        \textbf{AP-ID} & \textbf{Verantwortlich} & \textbf{Arbeitspaket} \\
        \midrule
        AP 3.1 & Drazen & JWT-Authentifizierung implementieren \\
        AP 3.2 & Drazen & Google OAuth Integration \\
        AP 3.3 & Drazen & Registrierung \& Login (Backend) \\
        AP 3.4 & Ereny  & Login- \& Registrierungsseite (Frontend) \\
        AP 3.5 & Ereny  & Profilverwaltung (Frontend) \\
        \bottomrule
    \end{tabularx}
    \label{tab:ap3}
\end{table}

\subsection{Phase 4: Kernfunktionen (Jänner 2026)}

\begin{table}[H]
    \centering
    \caption{Arbeitspakete Phase 4}
    \begin{tabularx}{\textwidth}{llX}
        \toprule
        \textbf{AP-ID} & \textbf{Verantwortlich} & \textbf{Arbeitspaket} \\
        \midrule
        AP 4.1 & Drazen & Task-CRUD-API (Backend) \\
        AP 4.2 & Drazen & Subtask-Verwaltung (Backend) \\
        AP 4.3 & Ereny  & Task-Ansicht \& Verwaltung (Frontend) \\
        AP 4.4 & Ereny  & Mitgliederverwaltung (Frontend \& Backend) \\
        AP 4.5 & Gemeinsam & Benachrichtigungssystem \\
        \bottomrule
    \end{tabularx}
    \label{tab:ap4}
\end{table}

\subsection{Phase 5: Kommunikation \& Echtzeit (Februar 2026)}

\begin{table}[H]
    \centering
    \caption{Arbeitspakete Phase 5}
    \begin{tabularx}{\textwidth}{llX}
        \toprule
        \textbf{AP-ID} & \textbf{Verantwortlich} & \textbf{Arbeitspaket} \\
        \midrule
        AP 5.1 & Drazen & WebSocket-Integration (Socket.io, Backend) \\
        AP 5.2 & Ereny  & Live-Chat-Komponente (Frontend) \\
        AP 5.3 & Drazen & Chatverlauf-Persistenz \& DSGVO-Löschkonzept \\
        AP 5.4 & Ereny  & Globale Suche (Frontend \& Backend) \\
        \bottomrule
    \end{tabularx}
    \label{tab:ap5}
\end{table}

\subsection{Phase 6: Kalender \& Dashboard (März 2026)}

\begin{table}[H]
    \centering
    \caption{Arbeitspakete Phase 6}
    \begin{tabularx}{\textwidth}{llX}
        \toprule
        \textbf{AP-ID} & \textbf{Verantwortlich} & \textbf{Arbeitspaket} \\
        \midrule
        AP 6.1 & Ereny  & Kalenderansicht (Frontend) \\
        AP 6.2 & Drazen & Kalender-API (Backend) \\
        AP 6.3 & Ereny  & Dashboard-Metriken Vollausbau:
                          Statusverteilung, Fortschritt, Trends \\
        AP 6.4 & Drazen & Backend-APIs für Datenvisualisierung \\
        \bottomrule
    \end{tabularx}
    \label{tab:ap6}
\end{table}

\subsection{Phase 7: Sicherheit \& Compliance (März--April 2026)}

\begin{table}[H]
    \centering
    \caption{Arbeitspakete Phase 7}
    \begin{tabularx}{\textwidth}{llX}
        \toprule
        \textbf{AP-ID} & \textbf{Verantwortlich} & \textbf{Arbeitspaket} \\
        \midrule
        AP 7.1 & Drazen & Zwei-Faktor-Authentifizierung (2FA/TOTP) \\
        AP 7.2 & Ereny  & Rollen-/Rechteprüfung serverseitig \&
                          Audit-Log \\
        AP 7.3 & Drazen & Sicherheitsmaßnahmen (Interceptors,
                          Middleware, Rate Limiting) \\
        AP 7.4 & Drazen & Secrets Management \& HTTPS-Vorgaben \\
        AP 7.5 & Drazen & Hintergrundprozesse \& Cronjobs \\
        AP 7.6 & Drazen & Deployment, Backup- \&
                          Restore-Strategie, Monitoring \\
        \bottomrule
    \end{tabularx}
    \label{tab:ap7}
\end{table}

\authorsection{Projektmanagement-Methodik}{Dražen Plavšić, Ereny Micail}

Für die Planung und Durchführung des Projekts wurde ein hybrides
Vorgehensmodell gewählt, das Elemente aus dem \textbf{Wasserfallmodell}
und \textbf{agilen Methoden} kombiniert:

\begin{description}
    \item[Wasserfallmodell] Die übergeordnete Projektstruktur folgt
          einem sequenziellen Ablauf mit klar definierten Phasen
          (Analyse, Design, Implementierung, Test). Dies ermöglichte
          eine strukturierte Planung und klare Meilensteine.
    \item[Agile Elemente] Innerhalb der Phasen wurden kurze
          Entwicklungszyklen (Sprints) mit regelmäßigen Besprechungen
          mit dem Betreuer durchgeführt. Änderungen an Anforderungen
          konnten so flexibel eingearbeitet werden.
    \item[Git-Flow] Die Versionsverwaltung folgte einer strukturierten
          Git-Flow-Strategie mit einem \texttt{main}-Branch für stabile
          Releases, einem \texttt{development}-Branch für die laufende
          Entwicklung und Feature-Branches für einzelne Arbeitspakete
          (z.\,B.\ \texttt{feature/35-dashboard-metriken}).
    \item[Issue-Tracking] Alle Arbeitspakete wurden als GitHub-Issues
          erfasst und mit Verantwortlichen, Labels und Meilensteinen
          versehen. Dies ermöglichte eine transparente Nachverfolgung
          des Projektfortschritts.
\end{description}


\authorsection{Anforderungsanalyse}{Dražen Plavšić, Ereny Micail}

Die Anforderungsanalyse bildet die Grundlage für die gesamte Entwicklung
von EffSync. Sie definiert, welche Funktionen das System bereitstellen
muss (funktionale Anforderungen) und welche Qualitätseigenschaften es
erfüllen soll (nicht-funktionale Anforderungen). Die Anforderungen wurden
gemeinsam mit dem Betreuer \textbf{Simon Gunacker} in der Kickoff-Phase
erarbeitet und im Laufe des Projekts iterativ verfeinert.

\subsection{Stakeholder}

Folgende Stakeholder wurden im Rahmen der Anforderungsanalyse identifiziert:

\begin{table}[H]
    \centering
    \caption{Stakeholder-Übersicht}
    \begin{tabularx}{\textwidth}{lXl}
        \toprule
        \textbf{Stakeholder} & \textbf{Interesse / Erwartung}
          & \textbf{Priorität} \\
        \midrule
        Schülerinnen \& Schüler
          & Einfache Aufgabenverwaltung, Kalender, Chat
          & Hoch \\
        Lehrpersonen
          & Überblick über Klassen, Aufgaben und Fortschritt
          & Hoch \\
        Unternehmen (Teams)
          & Skalierbare Projektverwaltung, Sicherheit, DSGVO
          & Hoch \\
        Betreuer (Simon Gunacker)
          & Technische Qualität, Dokumentation, Git-Workflow
          & Mittel \\
        Entwickler (Plavsic, Micail)
          & Wartbarkeit, klare Architektur, Testbarkeit
          & Mittel \\
        \bottomrule
    \end{tabularx}
    \label{tab:stakeholder}
\end{table}

\subsection{Funktionale Anforderungen}
\label{sec:funktionale-anforderungen}

Funktionale Anforderungen beschreiben, \textit{was} das System tun soll.
Sie sind nach Themenbereichen gruppiert und mit einer eindeutigen ID
sowie einer Priorität (Muss / Soll / Kann) versehen.

\subsubsection{Benutzerverwaltung \& Authentifizierung}

\begin{table}[H]
    \centering
    \caption{Funktionale Anforderungen -- Benutzerverwaltung}
    \begin{tabularx}{\textwidth}{llXl}
        \toprule
        \textbf{ID} & \textbf{Prio} & \textbf{Anforderung}
          & \textbf{Status} \\
        \midrule
        FA-01 & Muss & Benutzer können sich mit E-Mail und Passwort
                       registrieren und anmelden.
                     & \checkmark \\
        FA-02 & Muss & Passwörter werden gehasht (bcrypt) gespeichert.
                     & \checkmark \\
        FA-03 & Muss & Authentifizierung erfolgt über \gls{JWT}-Token
                       (Access- \& Refresh-Token).
                     & \checkmark \\
        FA-04 & Soll & Anmeldung via Google \gls{OAuth} ist möglich.
                     & \checkmark \\
        FA-05 & Soll & Benutzer können ihr Profil (Name, Avatar,
                       Passwort) bearbeiten.
                     & \checkmark \\
        FA-06 & Kann & Zwei-Faktor-Authentifizierung (\gls{2FA})
                       via \gls{TOTP} (z.\,B.\ Google Authenticator)
                       ist verfügbar.
                     & \checkmark \\
        FA-07 & Muss & Passwort-Reset per E-Mail ist möglich.
                     & \checkmark \\
        \bottomrule
    \end{tabularx}
    \label{tab:fa-auth}
\end{table}

\subsubsection{Rollen \& Berechtigungen}

\begin{table}[H]
    \centering
    \caption{Funktionale Anforderungen -- Rollen \& Berechtigungen}
    \begin{tabularx}{\textwidth}{llXl}
        \toprule
        \textbf{ID} & \textbf{Prio} & \textbf{Anforderung}
          & \textbf{Status} \\
        \midrule
        FA-10 & Muss & Das System verfügt über ein rollenbasiertes
                       Berechtigungssystem (\gls{RBAC}).
                     & \checkmark \\
        FA-11 & Muss & Rollen: Superuser, Admin, Manager, Member,
                       Viewer (Unternehmen) bzw.\ Superuser, Admin,
                       Teacher, Student (Schule).
                     & \checkmark \\
        FA-12 & Muss & Jede Rolle hat klar definierte Berechtigungen
                       für CRUD-Operationen.
                     & \checkmark \\
        FA-13 & Soll & Berechtigungen werden serverseitig geprüft
                       (Guards in NestJS).
                     & \checkmark \\
        FA-14 & Soll & Kritische Operationen werden in einem
                       Audit-Log protokolliert.
                     & \checkmark \\
        \bottomrule
    \end{tabularx}
    \label{tab:fa-rbac}
\end{table}

\subsubsection{Aufgabenverwaltung}

\begin{table}[H]
    \centering
    \caption{Funktionale Anforderungen -- Aufgabenverwaltung}
    \begin{tabularx}{\textwidth}{llXl}
        \toprule
        \textbf{ID} & \textbf{Prio} & \textbf{Anforderung}
          & \textbf{Status} \\
        \midrule
        FA-20 & Muss & Tasks können erstellt, bearbeitet, gelöscht
                       und einem Status zugewiesen werden.
                     & \checkmark \\
        FA-21 & Muss & Tasks haben folgende Felder: Titel,
                       Beschreibung, Status, Priorität,
                       Fälligkeitsdatum, Verantwortliche.
                     & \checkmark \\
        FA-22 & Muss & Tasks können Subtasks enthalten.
                     & \checkmark \\
        FA-23 & Soll & Tasks können Mitgliedern explizit zugewiesen
                       werden (Task-Member-Verwaltung).
                     & \checkmark \\
        FA-24 & Soll & Tasks können nach Status, Priorität und
                       Fälligkeitsdatum gefiltert werden.
                     & \checkmark \\
        FA-25 & Kann & Tasks können mit Tags/Labels versehen werden.
                     & \checkmark \\
        \bottomrule
    \end{tabularx}
    \label{tab:fa-tasks}
\end{table}

\subsubsection{Dashboard \& Metriken}

\begin{table}[H]
    \centering
    \caption{Funktionale Anforderungen -- Dashboard}
    \begin{tabularx}{\textwidth}{llXl}
        \toprule
        \textbf{ID} & \textbf{Prio} & \textbf{Anforderung}
          & \textbf{Status} \\
        \midrule
        FA-30 & Muss & Das Dashboard zeigt eine Übersicht aller
                       aktiven Projekte und Tasks.
                     & \checkmark \\
        FA-31 & Soll & Metriken: Gesamtanzahl Tasks, offene Tasks,
                       abgeschlossene Tasks, überfällige Tasks.
                     & \checkmark \\
        FA-32 & Soll & Statusverteilung wird als Diagramm
                       (Donut-Chart) visualisiert.
                     & \checkmark \\
        FA-33 & Soll & Wöchentliche und monatliche Trends werden
                       als Liniendiagramm dargestellt.
                     & \checkmark \\
        FA-34 & Kann & Fortschrittsbalken für einzelne Projekte
                       werden angezeigt.
                     & \checkmark \\
        FA-35 & Kann & Top-5-Mitglieder nach abgeschlossenen
                       Tasks werden angezeigt.
                     & \checkmark \\
        \bottomrule
    \end{tabularx}
    \label{tab:fa-dashboard}
\end{table}

\subsubsection{Kommunikation \& Benachrichtigungen}

\begin{table}[H]
    \centering
    \caption{Funktionale Anforderungen -- Kommunikation}
    \begin{tabularx}{\textwidth}{llXl}
        \toprule
        \textbf{ID} & \textbf{Prio} & \textbf{Anforderung}
          & \textbf{Status} \\
        \midrule
        FA-40 & Soll & Ein Live-Chat ermöglicht Echtzeit-Kommunikation
                       via WebSockets (Socket.io).
                     & \checkmark \\
        FA-41 & Soll & Chatnachrichten werden persistent gespeichert
                       (DSGVO-konform mit Löschkonzept).
                     & \checkmark \\
        FA-42 & Muss & Benutzer erhalten In-App-Benachrichtigungen
                       bei relevanten Ereignissen (z.\,B.\ neue
                       Task-Zuweisung).
                     & \checkmark \\
        FA-43 & Soll & Benachrichtigungen können als gelesen
                       markiert werden.
                     & \checkmark \\
        FA-44 & Kann & E-Mail-Benachrichtigungen bei kritischen
                       Ereignissen (z.\,B.\ überfällige Tasks).
                     & \checkmark \\
        \bottomrule
    \end{tabularx}
    \label{tab:fa-komm}
\end{table}

\subsubsection{Kalender \& Suche}

\begin{table}[H]
    \centering
    \caption{Funktionale Anforderungen -- Kalender \& Suche}
    \begin{tabularx}{\textwidth}{llXl}
        \toprule
        \textbf{ID} & \textbf{Prio} & \textbf{Anforderung}
          & \textbf{Status} \\
        \midrule
        FA-50 & Soll & Eine Kalenderansicht zeigt Tasks nach
                       Fälligkeitsdatum an.
                     & \checkmark \\
        FA-51 & Soll & Der Kalender unterstützt Monats-, Wochen-
                       und Tagesansicht.
                     & \checkmark \\
        FA-52 & Soll & Eine globale Suche ermöglicht das Auffinden
                       von Tasks, Projekten und Benutzern.
                     & \checkmark \\
        FA-53 & Kann & Suchergebnisse werden in Kategorien
                       gruppiert dargestellt.
                     & \checkmark \\
        \bottomrule
    \end{tabularx}
    \label{tab:fa-kalender}
\end{table}

\subsection{Nicht-funktionale Anforderungen}
\label{sec:nichtfunktionale-anforderungen}

Nicht-funktionale Anforderungen beschreiben \textit{wie} das System
seine Aufgaben erfüllen soll. Sie betreffen Qualitätseigenschaften
wie Performance, Sicherheit und Wartbarkeit.

\begin{table}[H]
    \centering
    \caption{Nicht-funktionale Anforderungen}
    \begin{tabularx}{\textwidth}{llX}
        \toprule
        \textbf{ID} & \textbf{Kategorie} & \textbf{Anforderung} \\
        \midrule
        NFA-01 & Performance
          & API-Antwortzeiten sollen unter 200\,ms liegen
            (bei normaler Last). \\
        NFA-02 & Performance
          & Das Frontend soll innerhalb von 3 Sekunden vollständig
            geladen sein. \\
        NFA-03 & Skalierbarkeit
          & Die Anwendung soll mindestens 100 gleichzeitige
            Benutzer unterstützen. \\
        NFA-04 & Sicherheit
          & Alle Passwörter werden mit bcrypt (Kostenfaktor 12)
            gehasht. \\
        NFA-05 & Sicherheit
          & Die Kommunikation erfolgt ausschließlich über HTTPS. \\
        NFA-06 & Sicherheit
          & \gls{JWT}-Tokens haben eine begrenzte Gültigkeitsdauer
            (Access-Token: 15\,min, Refresh-Token: 7\,Tage). \\
        NFA-07 & Sicherheit
          & Rate Limiting verhindert Brute-Force-Angriffe
            auf Login-Endpunkte. \\
        NFA-08 & Datenschutz
          & Die Anwendung ist \gls{DSGVO}-konform (Löschkonzepte,
            Datensparsamkeit). \\
        NFA-09 & Wartbarkeit
          & Der Code folgt den SOLID-Prinzipien und ist modular
            aufgebaut. \\
        NFA-10 & Wartbarkeit
          & Alle API-Endpunkte sind mit Swagger dokumentiert. \\
        NFA-11 & Usability
          & Die Benutzeroberfläche ist responsiv und auf mobilen
            Geräten nutzbar. \\
        NFA-12 & Usability
          & Die Anwendung unterstützt Dark Mode und Light Mode. \\
        NFA-13 & Verfügbarkeit
          & Das System soll eine Verfügbarkeit von 99\,\% anstreben. \\
        NFA-14 & Testbarkeit
          & Kritische Backend-Module werden durch Unit-Tests
            abgedeckt. \\
        NFA-15 & Kompatibilität
          & Die Anwendung ist mit den aktuellen Versionen von
            Chrome, Firefox und Edge kompatibel. \\
        \bottomrule
    \end{tabularx}
    \label{tab:nfa}
\end{table}

\subsection{Use-Case-Diagramm}
\label{sec:use-cases}

Die folgende Tabelle beschreibt die wichtigsten Use-Cases des Systems
gegliedert nach Akteuren:

\begin{table}[H]
    \centering
    \caption{Use-Case-Übersicht}
    \begin{tabularx}{\textwidth}{llX}
        \toprule
        \textbf{Akteur} & \textbf{UC-ID} & \textbf{Use-Case} \\
        \midrule
        Alle Benutzer
          & UC-01 & Registrieren \& Anmelden \\
          & UC-02 & Profil bearbeiten \\
          & UC-03 & Dashboard einsehen \\
          & UC-04 & Aufgaben anzeigen \& filtern \\
          & UC-05 & Kalender einsehen \\
          & UC-06 & Live-Chat nutzen \\
          & UC-07 & Globale Suche verwenden \\
          & UC-08 & Benachrichtigungen verwalten \\
        \midrule
        Member / Student
          & UC-10 & Aufgaben erstellen \& bearbeiten \\
          & UC-11 & Subtasks verwalten \\
          & UC-12 & Aufgaben als erledigt markieren \\
          & UC-13 & Im Rahmen einer Aufgabe Datein hochladen \& Note einsehen\\
        \midrule
        Manager / Teacher
          & UC-20 & Projekte \& Klassen verwalten \\
          & UC-21 & Mitglieder zuweisen \\
          & UC-22 & Metriken \& Trends einsehen \\
          & UC-23 & abgegebene Tasks bewerten \& Kommentieren \\
        \midrule
        Admin
          & UC-30 & Benutzerverwaltung \\
          & UC-31 & Rollen vergeben \\
          & UC-32 & Audit-Log einsehen \\
        \midrule
        Superuser
          & UC-40 & Systemkonfiguration \\
          & UC-41 & Backup \& Restore \\
          & UC-42 & Vollzugriff auf alle Daten \\
        \bottomrule
    \end{tabularx}
    \label{tab:usecases}
\end{table}

\subsection{Detaillierter Use-Case: Aufgabe erstellen}

\begin{table}[H]
    \centering
    \caption{UC-10: Aufgabe erstellen}
    \begin{tabularx}{\textwidth}{lX}
        \toprule
        \textbf{Feld} & \textbf{Beschreibung} \\
        \midrule
        Use-Case-ID   & UC-10 \\
        Name          & Aufgabe erstellen \\
        Akteur        & Member, Manager, Admin, Superuser \\
        Vorbedingung  & Benutzer ist authentifiziert und Mitglied
                        eines Projekts \\
        Nachbedingung & Neue Aufgabe ist in der Datenbank gespeichert
                        und im Frontend sichtbar \\
        Normalablauf  & 1. Benutzer klickt auf \glqq{}Neue Aufgabe\grqq{} \\
                      & 2. Formular mit Titel, Beschreibung, Status,
                           Priorität, Fälligkeitsdatum öffnet sich, Zeit muss eingetragen werden \\
                      & 3. Benutzer füllt Pflichtfelder aus \\
                      & 4. Benutzer weist optional Mitglieder zu \\
                      & 5. Benutzer speichert die Aufgabe \\
                      & 6. Backend validiert Eingaben und speichert
                           in PostgreSQL \\
                      & 7. Zugewiesene Mitglieder erhalten eine
                           Benachrichtigung \\
        Ausnahmen     & Pflichtfelder fehlen $\rightarrow$ Validierungsfehler \\
                      & Benutzer hat keine Berechtigung $\rightarrow$
                        403 Forbidden \\
        \bottomrule
    \end{tabularx}
    \label{tab:uc10}
\end{table}

\subsection{Detaillierter Use-Case: Anmeldung mit 2FA}

\begin{table}[H]
    \centering
    \caption{UC-01b: Anmeldung mit Zwei-Faktor-Authentifizierung}
    \begin{tabularx}{\textwidth}{lX}
        \toprule
        \textbf{Feld} & \textbf{Beschreibung} \\
        \midrule
        Use-Case-ID   & UC-01b \\
        Name          & Anmeldung mit 2FA \\
        Akteur        & Alle Benutzer (sofern 2FA aktiviert) \\
        Vorbedingung  & Benutzer hat 2FA in den Profileinstellungen
                        aktiviert und eine Authenticator-App
                        eingerichtet \\
        Nachbedingung & Benutzer ist authentifiziert und erhält
                        gültige JWT-Tokens \\
        Normalablauf  & 1. Benutzer gibt E-Mail und Passwort ein \\
                      & 2. Backend prüft Credentials \\
                      & 3. Backend erkennt aktive 2FA und fordert
                           TOTP-Code an \\
                      & 4. Benutzer öffnet Authenticator-App und
                           gibt 6-stelligen Code ein \\
                      & 5. Backend validiert TOTP-Code \\
                      & 6. Access- und Refresh-Token werden ausgestellt \\
        Ausnahmen     & Falscher TOTP-Code $\rightarrow$ 401 Unauthorized \\
                      & Code abgelaufen $\rightarrow$ Neuen Code anfordern \\
        \bottomrule
    \end{tabularx}
    \label{tab:uc01b}
\end{table}